% Problema 23 del Cálculo de Spivak
% Demostración de acotación para el cociente x/y

\section*{Problema 23 - Spivak: Demostración de acotación para $\left|\frac{x}{y} - \frac{x_0}{y_0}\right|$}
\addcontentsline{toc}{section}{Problema 23 - Spivak: Demostración de acotación para $\left|\frac{x}{y} - \frac{x_0}{y_0}\right|$}

\subsection*{Enunciado}

Sustituir los interrogantes del siguiente enunciado por expresiones que encierren $\varepsilon$, $x$ e $y$, de tal manera que la conclusión sea válida:

\vspace{0.3cm}

Si $y_0 \neq 0$ y 
\[
|y - y_0| < ? \quad \text{y} \quad |x - x_0| < ?,
\]
entonces $y \neq 0$ y
\[
\left|\frac{x}{y} - \frac{x_0}{y_0}\right| < \varepsilon.
\]

\subsection*{Observación Crucial}

\textbf{El problema es trivial en el sentido siguiente:} La solución se puede derivar sin casi ningún trabajo de los Problemas 21 y 22, usando la \textbf{reescritura algebraica}:
\[
\frac{x}{y} = x \cdot \frac{1}{y}.
\]

Por tanto:
\[
\frac{x}{y} - \frac{x_0}{y_0} = x \cdot \frac{1}{y} - x_0 \cdot \frac{1}{y_0}.
\]

Este es un \textit{producto} (aplicable el Problema 21) de dos diferencias: $|x - x_0|$ y $\left|\dfrac{1}{y} - \dfrac{1}{y_0}\right|$ (aplicable el Problema 22).

\subsection*{Solución: Valores para los interrogantes}

Los interrogantes deben reemplazarse por:

\[
\boxed{|y - y_0| < \min\left(\frac{|y_0|}{2}, \frac{\varepsilon|y_0|^2}{2(|x_0|+1)}\right)}
\]

\[
\boxed{|x - x_0| < \frac{\varepsilon}{2(|y_0^{-1}|+1)} = \frac{\varepsilon}{2(|y_0|^{-1}+1)} = \frac{\varepsilon|y_0|}{2(1+|y_0|)}}
\]

O más simplemente (eligiendo ambos de forma simétrica):

\[
\boxed{|y - y_0| < \min\left(\frac{|y_0|}{2}, \frac{\varepsilon|y_0|^2}{4}\right) \quad \text{y} \quad |x - x_0| < \frac{\varepsilon}{4(|y_0|+1)}}
\]

\subsection*{Demostración Detallada}

\begin{proof}

\textbf{Paso 1: Aplicar el Problema 22 (recíproca).}

Por el Problema 22, si 
\[
|y - y_0| < \min\left(\frac{|y_0|}{2}, \frac{\varepsilon'|y_0|^2}{2}\right)
\]
para cierto $\varepsilon' > 0$, entonces $y \neq 0$ y 
\[
\left|\frac{1}{y} - \frac{1}{y_0}\right| < \varepsilon'.
\]

Elegimos $\varepsilon' = \dfrac{\varepsilon}{2(|x_0|+1)}$, así que necesitamos:
\[
|y - y_0| < \min\left(\frac{|y_0|}{2}, \frac{\varepsilon|y_0|^2}{4(|x_0|+1)}\right).
\]

Esto garantiza que $y \neq 0$ y
\[
\left|\frac{1}{y} - \frac{1}{y_0}\right| < \frac{\varepsilon}{2(|x_0|+1)}.
\]

\textbf{Paso 2: Aplicar el Problema 21 (producto).}

Reescribimos:
\[
\frac{x}{y} - \frac{x_0}{y_0} = x \cdot \frac{1}{y} - x_0 \cdot \frac{1}{y_0}.
\]

Aplicamos el Problema 21 con:
\begin{itemize}
\item $x$ en lugar de $x$ (numerador izquierdo),
\item $\dfrac{1}{y}$ en lugar de $y$ (denominador, que es una función de $y$),
\item $x_0$ en lugar de $x_0$,
\item $\dfrac{1}{y_0}$ en lugar de $y_0$.
\end{itemize}

El Problema 21 requiere:
\begin{align*}
|x - x_0| &< \min\left(\frac{\varepsilon}{2(|y_0^{-1}|+1)}, 1\right) = \min\left(\frac{\varepsilon|y_0|}{2(1+|y_0|)}, 1\right), \\
\left|\frac{1}{y} - \frac{1}{y_0}\right| &< \frac{\varepsilon}{2(|x_0|+1)}.
\end{align*}

La segunda condición se satisface del Paso 1. Para la primera, basta:
\[
|x - x_0| < \frac{\varepsilon|y_0|}{2(1+|y_0|)}.
\]

Entonces, por el Problema 21:
\[
\left|x \cdot \frac{1}{y} - x_0 \cdot \frac{1}{y_0}\right| < \varepsilon.
\]

Por tanto:
\[
\left|\frac{x}{y} - \frac{x_0}{y_0}\right| < \varepsilon. \quad \checkmark
\]

\textbf{Resumen de condiciones:}

\begin{enumerate}
\item $|y - y_0| < \min\left(\dfrac{|y_0|}{2}, \dfrac{\varepsilon|y_0|^2}{4(|x_0|+1)}\right)$
\item $|x - x_0| < \dfrac{\varepsilon|y_0|}{2(1+|y_0|)}$
\end{enumerate}

\end{proof}

\subsection*{Simplificación e Interpretación}

\textbf{Versión simplificada:} Para propósitos prácticos, podemos escribir:

\[
\boxed{|y - y_0| < \frac{|y_0|}{2} \quad \text{y} \quad |x - x_0| < \frac{\varepsilon|y_0|^2}{4}}
\]

con la condición adicional (del Problema 22) de que:
\[
|y - y_0| < \frac{\varepsilon|y_0|^2}{4}.
\]

Es decir, ambas desigualdades deben cumplirse simultáneamente.

\textbf{¿Por qué el cociente es más restrictivo que el producto?}

\begin{itemize}
\item En el Problema 21 (producto $xy$), la condición es $|x - x_0| < \min\left(\dfrac{\varepsilon}{2(|y_0|+1)}, 1\right)$, que depende linealmente de $\varepsilon$.

\item En el Problema 23 (cociente $x/y$), la condición es $|x - x_0| < \dfrac{\varepsilon|y_0|}{2(1+|y_0|)}$, que también depende linealmente de $\varepsilon$.

\item Sin embargo, para $y$, la condición en el Problema 22 (recíproca) es $|y - y_0| < \min\left(\dfrac{|y_0|}{2}, \dfrac{\varepsilon|y_0|^2}{2}\right)$, que depende de $\varepsilon^2$ en la segunda parte.

\item Esto refleja que el cociente es más sensible a cambios en $y$ (el denominador) que en $x$ (el numerador).
\end{itemize}

\subsection*{Por qué el Problema es ``Trivial''}

El término ``trivial'' en el enunciado original se refiere a que **no se necesita nueva técnica**. Simplemente:

\begin{enumerate}
\item Reescribir $\dfrac{x}{y} = x \cdot y^{-1}$.
\item Aplicar Problema 21 (producto).
\item Aplicar Problema 22 (recíproca $y^{-1}$).
\item Combinar las dos hipótesis.
\end{enumerate}

En contraste, los Problemas 21 y 22 requieren análisis más profundos (acotamiento de multiplicadores, garantizar $y \neq 0$, etc.).

\subsection*{Continuidad del Cociente de Funciones}

Este resultado demuestra que si:
\begin{itemize}
\item $x \to x_0$ (función continua en $x_0$),
\item $y \to y_0$ con $y_0 \neq 0$ (función continua en $x_0$, no nula),
\end{itemize}
entonces:
\[
\frac{x}{y} \to \frac{x_0}{y_0}.
\]

Es decir, \textbf{el cociente de funciones continuas es continuo en puntos donde el denominador no es cero.}

\subsection*{Estructura Jerárquica de Problemas}

\begin{center}
\begin{tabular}{|c|c|c|c|}
\hline
\textbf{Problema} & \textbf{Operación} & \textbf{Usa} & \textbf{Dificultad} \\
\hline
20 & Suma/Resta & Desig. triangular & Baja \\
\hline
21 & Producto & Acotamiento de factores & Media \\
\hline
22 & Recíproca & Análisis de signos, $\min$ & Media-Alta \\
\hline
23 & Cociente & Problemas 21 + 22 & Alta (por síntesis) \\
\hline
24 & Suma asociativa & Inducción & Media \\
\hline
\end{tabular}
\end{center}

El Problema 23 está ``alto'' en la jerarquía aunque sea ``trivial'' porque \textbf{sintetiza} técnicas de problemas anteriores.

\subsection*{Ejemplos Numéricos}

\begin{ejemplo}
Sea $x_0 = 2$, $y_0 = 1$, $\varepsilon = 0.1$.

\textbf{Condiciones del Problema 23:}
\begin{align*}
|y - y_0| &< \min\left(\frac{1}{2}, \frac{0.1 \cdot 1^2}{4(2+1)}\right) = \min(0.5, 1/120) = 1/120 \approx 0.0083, \\
|x - x_0| &< \frac{0.1 \cdot 1}{2(1+1)} = \frac{0.1}{4} = 0.025.
\end{align*}

Si $|x - 2| < 0.025$ y $|y - 1| < 0.0083$, entonces $y \neq 0$ y
\[
\left|\frac{x}{y} - \frac{2}{1}\right| < 0.1.
\]

Es decir, si $x \in (1.975, 2.025)$ y $y \in (0.9917, 1.0083)$, entonces $\dfrac{x}{y}$ está dentro de $0.1$ de $2$. ✓
\end{ejemplo}

\begin{ejemplo}
Sea $x_0 = 1$, $y_0 = 0.5$, $\varepsilon = 0.05$.

\textbf{Condiciones:}
\begin{align*}
|y - y_0| &< \min\left(\frac{0.5}{2}, \frac{0.05 \cdot (0.5)^2}{4(1+1)}\right) = \min(0.25, 0.003125) = 0.003125, \\
|x - x_0| &< \frac{0.05 \cdot 0.5}{2(1+0.5)} = \frac{0.025}{3} \approx 0.0083.
\end{align*}

Si estas condiciones se satisfacen, entonces $\left|\dfrac{x}{y} - 2\right| < 0.05$. ✓
\end{ejemplo}

\subsection*{Observación Final}

El Problema 23 cierra el ciclo de operaciones aritméticas básicas (suma, resta, producto, cociente) desde el punto de vista de continuidad. Todos se heredan de dos propiedades fundamentales:
\begin{enumerate}
\item La \textbf{desigualdad triangular} (para suma/resta).
\item La \textbf{acotación de factores} (para producto y cociente).
\end{enumerate}

Estos dos pilares permiten construir teoría de continuidad para todas las funciones racionales y algebraicas.
